% ============================================================
%  §4.4 Αναγνώριση παραμέτρων μοντέλου
% ============================================================

\label{subsec:parameter_id}

Το θερμικό μοντέλο περιέχει 7 παραμέτρους ανά τροχό (διάνυσμα $\mathbf{p} \in
\mathbb{R}^7$, βλ. \cref{eq:param_vector}). Αυτές οι παράμετροι δεν είναι άμεσα μετρήσιμες
από τη γεωμετρία ή τις φυσικές ιδιότητες υλικού -- πρέπει να αναγνωριστούν από πειραματικά
δεδομένα. Η διαδικασία αυτή, γνωστή ως \emph{παραμετρική αναγνώριση} ή system
identification, αποτελεί ένα πρόβλημα βελτιστοποίησης: βρες εκείνο το $\mathbf{p}$ που
ελαχιστοποιεί τη διαφορά μεταξύ προσομοιωμένης και μετρημένης εξέλιξης θερμοκρασίας.

\subsubsection{Διατύπωση του προβλήματος}

Δεδομένης σειράς μετρήσεων θερμοκρασίας $T_{\rm meas}(s_k)$ (από αισθητήρες infrared ή
contact), το πρόβλημα ελαχιστοποίησης γράφεται:

\begin{equation}
    \hat{\mathbf{p}} = \operatorname{arg\,min}_{\mathbf{p} \in \mathcal{P}} \sum_{k=1}^{N}
    \Bigl[ T_{\rm sim}(s_k;\,\mathbf{p}) - T_{\rm meas}(s_k) \Bigr]^2
    \label{eq:param_id_obj}
\end{equation}

όπου $T_{\rm sim}(s_k;\,\mathbf{p})$ είναι το αποτέλεσμα ολοκλήρωσης της
\cref{eq:dTds_full} με παράμετρο $\mathbf{p}$, και $\mathcal{P}$ το σύνολο φυσικά
αποδεκτών τιμών (π.χ. $p_i > 0$, $c_w \in [0.1, 10]$).

Η βελτιστοποίηση εκτελείται με \texttt{lsqnonlin} (MATLAB), το οποίο υλοποιεί μέθοδο
trust-region-reflective για nonlinear least squares. Ο αλγόριθμος απαιτεί μόνο
αξιολογήσεις της συνάρτησης σφάλματος (χωρίς αναλυτική παράγωγο), εκτιμώντας το Jacobian
με πεπερασμένες διαφορές.

\subsubsection{Ευαισθησία και αναγνωρισιμότητα}

Ένα κλασικό πρόβλημα στην παραμετρική αναγνώριση ΟΔΕ είναι η \emph{αναγνωρισιμότητα}
(identifiability): μπορούν όλες οι παράμετροι να προσδιοριστούν μοναδικά από τα
παρατηρούμενα δεδομένα, ή υπάρχουν συνδυασμοί παραμέτρων που παράγουν ταυτόσημα
αποτελέσματα;

Στη συγκεκριμένη περίπτωση, η παράμετρος $c_w$ (scaling της θερμικής μάζας) εμφανίζεται
πάντα ως γινόμενο με $m_c c_{p,c}$ στον παρονομαστή -- αν και οι τρεις αυτές ποσότητες δεν
αναγνωρίζονται ξεχωριστά, η παράμετρος $c_w$ απορροφά τη συνολική αβεβαιότητα. Παρόμοιες
συσχετίσεις μπορεί να εμφανίζονται μεταξύ όρων πηγής και απώλειας.

% TODO: Ανάλυση ευαισθησίας (sensitivity analysis) για να εντοπίσεις ποιες παράμετροι
% επηρεάζουν περισσότερο την έξοδο σε ποιο τμήμα πίστας.
% Βοηθάει στην κατανόηση του μοντέλου και εντοπίζει αδύναμα σημεία δεδομένων.

\subsubsection{Αρχικές τιμές και bounds}

Η επιλογή αρχικών τιμών $\mathbf{p}_0$ είναι κρίσιμη λόγω της μη-κυρτότητας του
προβλήματος: ο \texttt{lsqnonlin} εγγυάται τοπικό ελάχιστο, όχι ολικό. Στην παρούσα
υλοποίηση χρησιμοποιούνται αρχικές τιμές βασισμένες σε τυπικές τιμές βιβλιογραφίας για
GT-class ελαστικά. Ένας multi-start προσέγγιση (διαφορετικά $\mathbf{p}_0$) μπορεί να
βελτιώσει την εμπιστοσύνη στο αποτέλεσμα.

% TODO: Δείξε αποτελέσματα αναγνώρισης: προσομοιωμένη vs. μετρημένη θερμοκρασία
% για τον γύρο αναγνώρισης (training lap) και για έναν άλλο γύρο (validation lap).
% Η διαφορά training/validation αξιολογεί γενίκευση του μοντέλου.
